%%==============================================================================
%% sections/04_method.tex
%%==============================================================================
\section{Method}
\label{sec:method}

\subsection{Synthetic-Data Curriculum}
\label{sec:method-data}
We construct a training corpus $\Data = \Data_{\mathrm{easy}} \cup
\Data_{\mathrm{med}} \cup \Data_{\mathrm{hard}}$ with increasing difficulty.
Difficulty is estimated by an ensemble of three verifier LLMs.

\subsection{Sparse-Attention Architecture}
\label{sec:method-arch}

\begin{figure}[t]
    \centering
    \begin{tikzpicture}
        \tikzset{
            stage/.style={rectangle, rounded corners=2pt, draw=bjtublue,
                          fill=bjtublue-light, minimum height=1.0cm,
                          minimum width=1.6cm, font=\sffamily\small},
            attn/.style ={rectangle, rounded corners=2pt, draw=bjtured,
                          fill=bjtured-light, minimum height=1.0cm,
                          minimum width=1.6cm, font=\sffamily\small},
            arrow/.style={-Latex, thick, bjtublue}
        }
        \node[stage] (emb) {Embed};
        \node[attn,  right=0.35cm of emb] (a1)  {Sparse Attn};
        \node[stage, right=0.35cm of a1]  (ff1) {FFN};
        \node[attn,  right=0.35cm of ff1] (a2)  {Sparse Attn};
        \node[stage, right=0.35cm of a2]  (ff2) {FFN};
        \node[stage, right=0.35cm of ff2] (head){LM Head};
        \foreach \a/\b in {emb/a1, a1/ff1, ff1/a2, a2/ff2, ff2/head}
            \draw[arrow] (\a) -- (\b);
        \node[below=0.2cm of ff1, font=\sffamily\footnotesize\color{ink-mute}]
            {$\times\,N$ layers};
    \end{tikzpicture}
    \caption{The BJTU-LLM decoder-only architecture with hardware-aware sparse
    attention (red blocks) interleaved with dense feed-forward layers.}
    \label{fig:arch}
\end{figure}

Sparse attention over a window partition $\mathcal{W}$ is defined as
\begin{equation}
    \mathrm{SparseAttn}(\mat{Q},\mat{K},\mat{V})
      = \bigoplus_{W\in\mathcal{W}}\softmax\!\left(\frac{\mat{Q}_W \mat{K}_W^\top}{\sqrt{d}}\right)\mat{V}_W.
    \label{eq:sparse-attn}
\end{equation}

\begin{theorem}{Complexity Bound}{thm:complexity}
Let $N$ denote sequence length and $d$ hidden dimension. The sparse-attention
operator in Eq.~\eqref{eq:sparse-attn} runs in time
$\mathcal{O}(N\sqrt{N}\cdot d)$, strictly less than the $\mathcal{O}(N^2 d)$
cost of dense attention for $N > 4$.
\end{theorem}

\begin{proof}
Each window contains $\sqrt{N}$ tokens; there are $\sqrt{N}$ windows; each
pair-wise interaction costs $\mathcal{O}(d)$. Summing yields
$\sqrt{N}\cdot\sqrt{N}\cdot\sqrt{N}\cdot d = N\sqrt{N}\,d$.
\end{proof}

\subsection{Two-Stage RLHF with Process Rewards}
\label{sec:method-rlhf}

Our training procedure alternates supervised fine-tuning (SFT) and RL updates
using Proximal Policy Optimization (PPO), as summarized in
Algorithm~\ref{alg:train}.

\begin{algorithm}[t]
\DontPrintSemicolon
\caption{Two-Stage RLHF Training for BJTU-LLM}
\label{alg:train}
\KwIn{Base model $\pmodel$, SFT dataset $\Data_{\mathrm{sft}}$,
      preference dataset $\Data_{\mathrm{pref}}$}
\KwOut{Aligned model $\pmodel^{\star}$}
\tcp*{Stage 1: Supervised fine-tuning}
\For{$t = 1,\dots,T_{\mathrm{sft}}$}{
    Sample mini-batch $(x, y) \sim \Data_{\mathrm{sft}}$\;
    $\params \leftarrow \params - \eta \nabla_{\params} \Loss_{\mathrm{CE}}(\params)$\;
}
\tcp*{Stage 2: Train process reward model $r_\phi$}
Fit $r_\phi$ on $\Data_{\mathrm{pref}}$ via ranking loss\;
\tcp*{Stage 3: PPO with process rewards}
\For{$t = 1,\dots,T_{\mathrm{rl}}$}{
    Sample prompt $x$, roll out $y \sim \pmodel(\cdot|x)$\;
    Compute step-wise rewards $\{r_\phi(x, y^{(\leq k)})\}_{k=1}^{T}$\;
    Update $\params$ via PPO clipping objective\;
}
\Return $\pmodel^{\star}$\;
\end{algorithm}
